%% ma: L10-B13-0076
%% lop: 10
%% chuong: 5
%% bai: 13
%% dang: 10.13.11
%% loai: TN4
%% muc: NB
%% dap_an: "B"
%% nguon: "chuyen-de-cac-so-dac-trung-cua-mau-so-lieu-khong-ghep-nhom-toan-10-chuong-trinh-moi.pdf (GV. Phan Nhật Linh) (P03), Bài 02 Dạng 2 Phần I Câu 1"
%% kien_thuc: "Lớp 10 Bài 13 (trung vị, tứ phân vị)"
%% trang_thai: chinh-thuc
%% ghi_chu: "Câu 16 gốc trùng hoàn toàn câu này, đã bỏ. Tự duyệt theo yêu cầu giáo viên 22/10/2026."
\begin{ex}
Điểm kiểm tra Toán của một nhóm học sinh như sau: $7$; $9$; $6$; $10$; $5$; $8$; $4$. Tứ phân vị thứ hai của mẫu số liệu là
\choice{$6$}{\True $7$}{$9$}{$5$}
\loigiai{Sắp xếp: $4$; $5$; $6$; $7$; $8$; $9$; $10$. $Q_2$ là trung vị: $Q_2=7$. Chọn B.}
\end{ex}
